% ==============================================================================
% บทที่ 3: วิธีการดำเนินโครงงาน (Research Methodology and System Implementation)
% โครงการ: การพัฒนาแพลตฟอร์มสำหรับการสนับสนุนการวินิจฉัยรอยโรคทางช่องปาก
%          และผสมผสานการใช้งานระบบผู้เชี่ยวชาญสำหรับทันตแพทย์
% ==============================================================================

\chapter{วิธีการดำเนินโครงงาน}
\label{chapter:methodology}

โครงงานนี้มีเป้าหมายในการพัฒนาแพลตฟอร์มเว็บแอปพลิเคชันเพื่อสนับสนุนการตัดสินใจทางคลินิกสำหรับทันตแพทย์ (Clinical Decision Support Platform) โดยมุ่งเน้นการนำแบบจำลองปัญญาประดิษฐ์ที่ได้รับการวิจัยและพัฒนาขึ้นจากโครงงานของรุ่นพี่ในปีก่อนหน้าทั้ง 2 กลุ่ม ได้แก่ แบบจำลองตรวจจับวัตถุ (YOLO) และแบบจำลองการเรียนรู้เชิงลึกแบบหลายโมดาลิตี (Multimodal Deep Learning: DenseNet-121 ร่วมกับ WangchanBERTa) นำมารวบรวมและบูรณาการเข้าด้วยกันภายใต้สถาปัตยกรรมระบบไมโครเซอร์วิส เพื่อให้ทันตแพทย์และบุคลากรทางการแพทย์สามารถเลือกใช้งานแบบจำลอง (Alternative Model Selection) ให้เหมาะสมกับบริบทและความต้องการในการตรวจคัดกรองรอยโรคแต่ละกรณี พร้อมทั้งพัฒนาส่วนติดต่อผู้ใช้งานที่ทันสมัยและสะดวกต่อการปฏิบัติงานจริง

ขั้นตอนและระเบียบวิธีวิจัยในการดำเนินโครงงานแบ่งออกเป็น 6 ส่วนหลัก ดังนี้:
\begin{enumerate}
    \item การรวบรวมและวิเคราะห์คุณลักษณะของแบบจำลองปัญญาประดิษฐ์จากงานวิจัยรุ่นพี่
    \item การออกแบบสถาปัตยกรรมระบบแบบไมโครเซอร์วิส 3 คอนเทนเนอร์ (System Architecture \& Microservices)
    \item การออกแบบผังกรณีการใช้งานและผังกิจกรรมการทำงานของระบบ (Use Case \& Activity Diagram)
    \item การพัฒนาส่วนติดต่อผู้ใช้งานสำหรับทันตแพทย์ (User Interface Design)
    \item การจัดเตรียมสภาพแวดล้อม การคอนเทนเนอร์ไรซ์ และการทดสอบระบบอัตโนมัติ
    \item สรุปกรอบการดำเนินงานเพื่อเตรียมการประเมินผลในบทที่ 4
\end{enumerate}

% ------------------------------------------------------------------------------
% 3.1 การรวบรวมและวิเคราะห์แบบจำลองปัญญาประดิษฐ์จากงานวิจัยรุ่นพี่
% ------------------------------------------------------------------------------
\section{การรวบรวมและวิเคราะห์แบบจำลองปัญญาประดิษฐ์จากงานวิจัยรุ่นพี่}
\label{sec:ai_model_consolidation}

ในการดำเนินงานขั้นแรก คณะผู้จัดทำได้นำไฟล์โมเดลที่ผ่านการฝึกฝน (Pre-trained Weights), ซอร์สโค้ดการอนุมานผล (Inference Pipeline), และชุดกฎเกณฑ์ทางคลินิกจากปริญญานิพนธ์ของรุ่นพี่ทั้งสองกลุ่ม มารวบรวมและวิเคราะห์โครงสร้างข้อมูลนำเข้าและส่งออก (Input/Output Interface) เพื่อเตรียมการผสานรวมเข้าสู่ระบบ:

\subsection{แบบจำลองตรวจจับวัตถุ YOLO จากรุ่นพี่กลุ่มที่ 1 (ปีการศึกษา 2567)}
แบบจำลองของรุ่นพี่กลุ่มที่ 1 (เปี่ยมเดช และ ปุณณวิช) \cite{piamdech2024oral} เป็นแบบจำลองในกลุ่ม Object Detection ได้แก่ สถาปัตยกรรม YOLOv7 (\texttt{yolov7.pt}) และ YOLOv8 (\texttt{yolov8.pt}) ซึ่งมีลักษณะการทำงานสำคัญดังนี้:
\begin{enumerate}
    \item \textbf{ข้อมูลนำเข้า (Input):} ภาพถ่ายรอยโรคในช่องปากที่ผ่านการแปลงเป็นปริภูมิสี BGR และปรับขนาดภาพด้วยวิธี Letterbox ขนาด $640 \times 640$ พิกเซล ควบคู่กับข้อความอาการของผู้ป่วย
    \item \textbf{การตรวจจับและคัดกรองกรอบรอยโรค:} แบบจำลองจะทำการสกัดคุณลักษณะและทำนายพิกัด Bounding Box ในรูปแบบ $[x, y, w, h]$ พร้อมคะแนนความเชื่อมั่น โดยตั้งเกณฑ์คัดกรองขั้นต่ำ (Confidence Threshold) ไว้ที่ 0.25 จากนั้นใช้กระบวนการคัดกรองกรอบรอยโรคที่ซ้อนทับกัน (Concentric Box \& Overlap Suppression) โดยหากกรอบของคลาสเดียวกันมีค่าความทับซ้อน (IoU) เกินร้อยละ 40 หรือมีพื้นที่ครอบคลุมอยู่ในอีกกรอบหนึ่งเกินร้อยละ 70 จะทำการตัดกรอบที่มีคะแนนน้อยกว่าออกไป
    \item \textbf{การจับคู่กฎเกณฑ์อาการทางคลินิก (Clinical Rule Matching):} นำข้อความอาการของผู้ป่วยมาตัดแยกคำสำคัญ (Tokens) เพื่อเปรียบเทียบกับพจนานุกรมอาการเฉพาะของรอยโรค (\texttt{PATIENT\_CONDITIONS}) เช่น หากตรวจพบรอยโรค Leukoplakia และพบคำสำคัญ ``สูบบุหรี่'', ``ดื่มเหล้า'', หรือ ``เคี้ยวหมาก'' จะถือว่าอาการตรงกัน
    \item \textbf{การคำนวณคะแนนความมั่นใจขั้นสุดท้าย:}
    \begin{itemize}
        \item กรณีอาการตรงกับรอยโรค: ปรับคะแนนเพิ่มตามสูตร $\text{Score} = \frac{\text{conf}}{2} + 0.50$
        \item กรณีอาการไม่ตรงกับรอยโรค: ปรับคะแนนลดลงตามสูตร $\text{Score} = \frac{\text{conf}}{2}$ (หากเดิมเกิน 0.60) หรือคงค่าเดิมไว้
    \end{itemize}
    \item \textbf{ผลลัพธ์ที่ได้ (Output):} ภาพถ่ายที่ผ่านการวาดกรอบ Bounding Box พร้อมป้ายกำกับชื่อโรค ค่าเปอร์เซ็นต์ความเชื่อมั่น และสีเส้นกรอบจำเพาะตามชนิดรอยโรค
\end{enumerate}

\subsection{แบบจำลองการเรียนรู้เชิงลึกแบบหลายโมดาลิตีจากรุ่นพี่กลุ่มที่ 2 (ปีการศึกษา 2568)}
แบบจำลองของรุ่นพี่กลุ่มที่ 2 (รตา และ วรดา) \cite{unnahasarakhon2025multimodal} พัฒนาขึ้นบนแนวคิด Multimodal Deep Learning ซึ่งผสานแบบจำลองภาพและข้อความเข้าด้วยกัน ประกอบด้วย:
\begin{enumerate}
    \item \textbf{โมดูลประมวลผลภาพ (DenseNet-121):} โหลดน้ำหนักที่ผ่านการเทรน (\texttt{best\_densenet121\_img.pth}) รับภาพขนาด $600 \times 600$ พิกเซล และทำ Normalization ด้วยค่าเฉลี่ยและส่วนเบี่ยงเบนมาตรฐานจำเพาะของชุดข้อมูลรอยโรคช่องปาก สกัดเวกเตอร์คุณลักษณะขนาด $1 \times 1024$ มิติ ผ่านชั้น Global Average Pooling (GAP)
    \item \textbf{โมดูลประมวลผลภาษาไทย (WangchanBERTa):} โหลดน้ำหนัก (\texttt{best\_text.pth}) รับข้อความอาการภาษาไทย เข้าสู่ SentencePiece Tokenizer เพื่อแปลงเป็น Input IDs และ Attention Mask แล้วส่งผ่าน Transformer Encoder เพื่อสกัดเวกเตอร์ตัวแทนบริบทจากโทเคน \texttt{[CLS]} ขนาด $1 \times 768$ มิติ
    \item \textbf{การรวมผลลัพธ์ระดับการตัดสินใจ (Late Decision Fusion):} นำคะแนนความน่าจะเป็นของโมเดลภาพ ($P_{img}$) และโมเดลข้อความ ($P_{text}$) มารวมกันด้วยการถ่วงน้ำหนักภาพร้อยละ 60 และข้อความร้อยละ 40:
    \begin{equation}
        P_{final} = 0.60 \times P_{img} + 0.40 \times P_{text}
        \label{eq:ch3_late_fusion}
    \end{equation}
    \item \textbf{การสร้างภาพความร้อนอธิบายผล (Grad-CAM Heatmap):} คำนวณเกรเดียนต์จากเลเยอร์สุดท้ายของ DenseNet-121 เพื่อสร้างแผนที่ความร้อนแสดงตำแหน่งรอยโรคที่ AI ใช้ในการตัดสินใจ และนำไปทับซ้อนลงบนภาพจริง
    \item \textbf{ผลลัพธ์ที่ได้ (Output):} ชื่อรอยโรคที่ทำนาย, เปอร์เซ็นต์ความเชื่อมั่นรวม, และภาพผลลัพธ์ที่ซ้อนทับด้วย Grad-CAM Heatmap
\end{enumerate}

\subsection{การเปรียบเทียบคุณลักษณะของแบบจำลองทั้งสองระบบ}
เพื่อให้ผู้ใช้งานเข้าใจความแตกต่างและสามารถเลือกใช้งานได้อย่างถูกต้อง ตารางที่ \ref{tab:model_comparison} แสดงการเปรียบเทียบจุดเด่นและคุณลักษณะเชิงเทคนิคของทั้งสองแบบจำลอง:

\begin{table}[htbp]
\centering
\small
\caption{การเปรียบเทียบคุณสมบัติเชิงเทคนิคระหว่างแบบจำลอง YOLO และแบบจำลอง Multimodal}
\label{tab:model_comparison}
\begin{tabular}{p{3.5cm}p{5.0cm}p{5.0cm}}
\toprule
\textbf{มิติการเปรียบเทียบ} & \textbf{แบบจำลอง YOLO (รุ่นพี่กลุ่ม 1)} & \textbf{แบบจำลอง Multimodal (รุ่นพี่กลุ่ม 2)} \\
\midrule
สถาปัตยกรรมหลัก & YOLOv7 / YOLOv8 & DenseNet-121 + WangchanBERTa \\
รูปแบบการทำงาน & การตรวจจับวัตถุ (Object Detection) & การจำแนกภาพแบบหลายโหมด (Multimodal) \\
การระบุตำแหน่งรอยโรค & \textbf{ระบุได้ด้วย Bounding Box ชัดเจน} & ระบุด้วยภาพความร้อน Grad-CAM Heatmap \\
การผสานข้อมูลประวัติ & ใช้กฎเงื่อนไขภายนอก (Rule-based) & \textbf{รวมด้วย Late Fusion ภายในโมเดล (60:40)} \\
ความแม่นยำทางคลินิก & ร้อยละ 80--86 & \textbf{ร้อยละ 95.5 (ทนทานต่อข้อมูลขัดแย้ง)} \\
การจัดการกรณีภาพไม่ชัด & เสี่ยงต่อภาวะ No Detection & ยังคงสามารถใช้ข้อมูลข้อความช่วยประเมินได้ \\
ความเหมาะสมในการใช้งาน & \textbf{เหมาะสำหรับการตรวจหาตำแหน่งแผล} & \textbf{เหมาะสำหรับการประเมินความเสี่ยงมะเร็ง} \\
\bottomrule
\end{tabular}
\end{table}

\subsection{ยุทธศาสตร์และเหตุผลในการรวมแบบจำลองเข้าด้วยกัน (Model Consolidation Strategy)}
จากการวิเคราะห์แบบจำลองทั้งสองกลุ่ม คณะผู้จัดทำพบว่าไม่มีแบบจำลองใดสมบูรณ์แบบในทุกมิติ โดยแต่ละแบบจำลองมีข้อได้เปรียบที่ส่งเสริมซึ่งกันและกัน (Complementary Strengths):
\begin{enumerate}
    \item \textbf{การเสริมจุดเด่นด้านตำแหน่งและบริบท (Localization vs Contextual Understanding):} แบบจำลอง YOLO มีจุดแข็งในการระบุพิกัดและขนาดของรอยโรคอย่างชัดเจนผ่าน Bounding Box ซึ่งจำเป็นต่อการวางแผนตัดชิ้นเนื้อ (Biopsy) ในขณะที่แบบจำลอง Multimodal มีความโดดเด่นในการประมวลผลข้อความอาการและพฤติกรรมเสี่ยงของผู้ป่วย เช่น การสูบบุหรี่จัด การดื่มสุรา หรืออาการเจ็บแสบ ทำให้ประเมินความเสี่ยงมะเร็งได้รอบด้านกว่า
    \item \textbf{ความยืดหยุ่นในการสนับสนุนการตัดสินใจทางคลินิก (Clinical Workflow Flexibility):} ในบริบทคลินิกจริง ผู้ป่วยแต่ละรายมีอาการและลักษณะรอยโรคที่แตกต่างกัน หากรอยโรคมีขอบเขตชัดเจน ทันตแพทย์อาจต้องการดูขอบเขตแผลด้วยโมเดล YOLO แต่หากรอยโรคมีลักษณะก้ำกึ่ง แบนราบ หรือมีความเสี่ยงแฝง ทันตแพทย์จะต้องการผลการประเมินจาก Multimodal เพื่อดูภาพความร้อน Grad-CAM และระดับความเสี่ยง
    \item \textbf{การตรวจสอบยืนยันข้ามแบบจำลองเพื่อความเห็นที่สอง (Cross-Model Validation / Second Opinion):} แพลตฟอร์มที่รวมทั้งสองแบบจำลองไว้ในระบบเดียวกัน เปิดโอกาสให้ทันตแพทย์สามารถส่งตรวจภาพและอาการเคสเดียวกันผ่านทั้งสองแบบจำลองเพื่อเปรียบเทียบผลลัพธ์ หากผลการวิเคราะห์สอดคล้องกัน จะช่วยเพิ่มความมั่นใจในการวินิจฉัยอย่างมีนัยสำคัญ
\end{enumerate}

ดังนั้น เป้าหมายหลักของโครงงานนี้จึงไม่ใช่การฝึกแบบจำลองใหม่ตั้งแต่ต้น แต่เป็นการสร้าง \textbf{Consolidated AI Platform} ที่นำแบบจำลองของรุ่นพี่ทั้งสองชุดมารวมเป็นระบบเดียว ออกแบบกลไกการสลับและเลือกใช้งานแบบจำลอง (Alternative Model Selection) และแปลงผลลัพธ์ให้อยู่ในมาตรฐานเดียวกัน (Unified Schema) เพื่อพร้อมสำหรับการใช้งานจริงในคลินิกทันตกรรม

% ------------------------------------------------------------------------------
% 3.2 การออกแบบสถาปัตยกรรมระบบและการบูรณาการแบบจำลอง
% ------------------------------------------------------------------------------
\section{การออกแบบสถาปัตยกรรมระบบและการบูรณาการแบบจำลอง}
\label{sec:system_architecture_design}

\subsection{สถาปัตยกรรมไมโครเซอร์วิส (Microservices Architecture)}
เพื่อแก้ไขปัญหาการผูกติดกันของระบบและลดการกินทรัพยากรของสถาปัตยกรรมแบบ Monolithic เดิม ระบบใหม่ได้รับการออกแบบตามสถาปัตยกรรม Microservices โดยแบ่งการทำงานออกเป็น 3 คอนเทนเนอร์หลัก (Three-Tier Microservices) ที่แยกหน้าที่กันอย่างอิสระและสื่อสารกันผ่านระบบเครือข่ายภายใน (\texttt{oralscan-network}):
\begin{enumerate}
    \item \textbf{Frontend Web Service (\texttt{oralscan-frontend}, พอร์ต 5173):} ส่วนติดต่อผู้ใช้งานฝั่งหน้าบ้าน พัฒนาด้วย React ร่วมกับ Vite และ Tailwind CSS ออกแบบในรูปแบบ Single Page Application (SPA) เพื่อให้ทันตแพทย์สามารถใช้งานได้อย่างสะดวก รวดเร็ว รองรับการอัปโหลดรูปภาพแบบ Drag-and-Drop มีระบบเลือกโมเดล (Model Selector) และเลือกอาการทางคลินิก (Symptoms Checkbox) โดยสื่อสารกับหลังบ้านผ่าน RESTful API
    \item \textbf{Backend API Gateway Service (\texttt{oralscan-backend}, พอร์ต 8000):} ส่วนหลังบ้านที่พัฒนาด้วย FastAPI ทำหน้าที่เป็นจุดรับคำขอภายนอกเพียงจุดเดียว (Single Entry Point) คอยตรวจสอบความถูกต้องของข้อมูลนำเข้า (Data Validation) ผ่าน Pydantic Schema, จัดการระบบรักษาความปลอดภัยและการจำกัดสิทธิ์ (CORS), จัดเก็บและให้บริการไฟล์รูปภาพทั้งภาพต้นฉบับและภาพผลลัพธ์ใน Static Directory (\texttt{/app/static}), จัดการเส้นทางคำขอ (Dynamic Model Router), และแปลงผลลัพธ์จาก AI ให้อยู่ในมาตรฐานเดียวกัน (Unified Response) พร้อมประเมินระดับความเสี่ยงมะเร็งทางคลินิก
    \item \textbf{AI Inference Engine Service (\texttt{oralscan-ai-engine}, พอร์ต 8001):} เซอร์วิสประมวลผลโมเดลปัญญาประดิษฐ์ความเร็วสูง บรรจุสภาพแวดล้อมทางคณิตศาสตร์และแมชชีนเลิร์นนิง ได้แก่ PyTorch, Ultralytics YOLO, OpenCV-headless และ HuggingFace Transformers โดยใช้ Thread-Safe Singleton Pattern ในการโหลดน้ำหนักของโมเดลรุ่นพี่ทั้งสองกลุ่มขึ้นสู่หน่วยความจำ (ได้แก่ \texttt{YoloService} และ \texttt{MultimodalService}) เพื่อป้องกันปัญหาโหลดโมเดลซ้ำซ้อน ประหยัด RAM และรองรับการประมวลผลคำขอพร้อมกันอย่างมีเสถียรภาพ
\end{enumerate}

\subsection{รูปแบบการเลือกใช้งานแบบจำลอง (Alternative Model Selection)}
เมื่อมีคำขอส่งตรวจเข้ามาผ่านเอนด์พอยต์ \texttt{/api/assess} ผู้ใช้งานสามารถส่งค่าพารามิเตอร์ \texttt{modelType} เพื่อกำหนดเส้นทางการประมวลผล (Dynamic Model Routing) ได้ตามความต้องการทางคลินิก ดังแสดงในรูปที่ \ref{fig:model_pipeline_diagram}:
\begin{itemize}
    \item \textbf{กรณีเลือกโมเดล YOLO (หรือ YOLOv7 / YOLOv8):} API Gateway จะส่งงานต่อไปยังฟังก์ชัน \texttt{infer\_yolo} ภายใน AI Engine โดยนำภาพรอยโรคเข้าสู่โมเดลตรวจจับวัตถุเพื่อค้นหา Bounding Box พร้อมทั้งนำอาการที่ผู้ใช้กรอกเข้าสู่กลไกจับคู่กฎเกณฑ์ทางคลินิก (Clinical Rule Matching) เพื่อคำนวณคะแนนความเชื่อมั่นสุทธิ และวาดกรอบสีจำเพาะลงบนภาพ
    \item \textbf{กรณีเลือกโมเดล DenseNet-BERT (Multimodal):} API Gateway จะส่งงานไปยังฟังก์ชัน \texttt{infer\_densenet\_bert} ภายใน AI Engine โดยแยกการประมวลผลภาพผ่าน DenseNet-121 และข้อความอาการผ่าน WangchanBERTa จากนั้นรวมผลด้วยกลไก Late Decision Fusion (น้ำหนักภาพร้อยละ 60 และข้อความร้อยละ 40) พร้อมสร้างภาพความร้อน Grad-CAM Heatmap
\end{itemize}

% แผนภาพขั้นตอนการรวมและการเลือกแบบจำลอง AI (เว้นกรอบ Placeholder สำหรับใส่ภาพกราฟิก)
\begin{figure}[htbp]
\centering
\fbox{
    \parbox{0.92\textwidth}{
        \centering
        \vspace{2.5cm}
        \textbf{[พื้นที่สำหรับใส่ภาพ: แผนภาพขั้นตอนการบูรณาการและการเลือกแบบจำลอง AI (AI Model Consolidation \& Selection Pipeline)]}\\[0.3cm]
        {\footnotesize \textit{(Inference Pipeline: Image \& Symptoms $\rightarrow$ Model Selection Router $\rightarrow$ [Branch A: YOLOv7/v8 + Rule Matching] vs [Branch B: Multimodal DenseNet-121 + WangchanBERTa] $\rightarrow$ Unified Schema Output)}}\\[0.2cm]
        {\scriptsize สามารถใส่ภาพ \texttt{images/model\_selection\_pipeline.png} ลงในจุดนี้}
        \vspace{2.5cm}
    }
}
\caption{แผนภาพขั้นตอนการบูรณาการและการเลือกแบบจำลองปัญญาประดิษฐ์ (AI Model Consolidation and Selection Pipeline)}
\label{fig:model_pipeline_diagram}
\end{figure}

\subsection{การกำหนดรูปแบบข้อมูลผลลัพธ์มาตรฐาน (Unified Assessment Schema)}
ไม่ว่าผู้ใช้งานจะเลือกประมวลผลด้วยแบบจำลองใด API Gateway จะส่งคืนผลลัพธ์ในโครงสร้าง JSON มาตรฐานเดียวกันเสมอ (\texttt{UnifiedAssessmentResponse}) ซึ่งประกอบด้วย:
\begin{itemize}
    \item \texttt{model\_used}: สถาปัตยกรรมโมเดลที่ใช้งานจริงในการตรวจ
    \item \texttt{disease\_name}: ชื่อรอยโรคเป้าหมายที่ตรวจพบ (Ulcer, Candidiasis, Linea alba buccalis, Leukoplakia, หรือ Lichen planus)
    \item \texttt{risk\_level}: ระดับความเสี่ยงต่อมะเร็งช่องปาก (Low, Medium, High)
    \item \texttt{confidence\_score}: ค่าคะแนนความเชื่อมั่นของแบบจำลอง (ทศนิยม 0.00 ถึง 1.00)
    \item \texttt{bounding\_boxes}: พิกัดกรอบรอยโรค (มีข้อมูลเมื่อใช้โมเดล YOLO และเป็นค่าว่างเมื่อใช้ DenseNet-BERT)
    \item \texttt{processed\_image\_url}: URL สำหรับดึงภาพที่ผ่านการวาดกรอบหรือ Grad-CAM ไปแสดงผล
\end{itemize}

\subsection{เกณฑ์การประเมินระดับความเสี่ยงมะเร็งทางคลินิก (Clinical Risk Level)}
เพื่อช่วยให้ทันตแพทย์จัดลำดับความเร่งด่วนในการรักษาได้ทันที ระบบได้กำหนดฟังก์ชันประเมินระดับความเสี่ยงมะเร็ง $\text{RiskLevel}(D, C)$ โดยพิจารณาจากชนิดรอยโรค ($D$) และคะแนนความเชื่อมั่น ($C$):
\begin{equation}
    \text{RiskLevel}(D, C) = 
    \begin{cases}
        \text{High (สูง)} & \text{เมื่อ } D \in \{\text{Leukoplakia}, \text{Ulcer}\} \text{ และ } C \ge 0.40 \\
        \text{Medium (ปานกลาง)} & \text{เมื่อ } D \in \{\text{Leukoplakia}, \text{Ulcer}\} \text{ และ } C < 0.40 \\
        \text{Medium (ปานกลาง)} & \text{เมื่อ } D \in \{\text{Lichen Planus}\} \\
        \text{Low (ต่ำ)} & \text{เมื่อ } D \in \{\text{Candidiasis}, \text{Linea Alba}\} \\
        \text{Low (ต่ำ)} & \text{เมื่อ } D \in \{\text{No Detection}, \text{Normal}\}
    \end{cases}
    \label{eq:ch3_risk_level}
\end{equation}

% แผนภาพสถาปัตยกรรมระบบ (เว้นกรอบ Placeholder สำหรับใส่ภาพกราฟิก)
\begin{figure}[htbp]
\centering
\fbox{
    \parbox{0.92\textwidth}{
        \centering
        \vspace{2.5cm}
        \textbf{[พื้นที่สำหรับใส่ภาพ: แผนภาพสถาปัตยกรรมระบบไมโครเซอร์วิส 3 คอนเทนเนอร์]}\\[0.3cm]
        {\footnotesize \textit{(System Architecture Diagram: Frontend Service (Port 5173) $\leftrightarrow$ Backend API Gateway (Port 8000) $\leftrightarrow$ AI Inference Engine (Port 8001))}}\\[0.2cm]
        {\scriptsize สามารถใส่ภาพ \texttt{images/system\_architecture.png} หรือรูปผังระบบที่วาดไว้ลงในจุดนี้}
        \vspace{2.5cm}
    }
}
\caption{แผนภาพสถาปัตยกรรมระบบไมโครเซอร์วิส 3 คอนเทนเนอร์ (Frontend, Backend API Gateway และ AI Inference Engine)}
\label{fig:system_architecture_diagram}
\end{figure}

% ------------------------------------------------------------------------------
% 3.3 การออกแบบผังกรณีการใช้งานและผังกิจกรรมการทำงานของระบบ
% ------------------------------------------------------------------------------
\section{การออกแบบผังกรณีการใช้งานและผังกิจกรรมการทำงานของระบบ}
\label{sec:usecase_and_workflow}

\subsection{ผังกรณีการใช้งานของระบบ (Use Case Diagram)}
การออกแบบผังกรณีการใช้งาน (Use Case Diagram) มีวัตถุประสงค์เพื่อกำหนดขอบเขตหน้าที่และความสามารถของระบบสนับสนุนการวินิจฉัยรอยโรคทางช่องปาก โดยแบ่งกลุ่มผู้มีส่วนร่วมกับระบบ (Actors) ออกเป็น 2 กลุ่มหลัก ได้แก่:
\begin{enumerate}
    \item \textbf{ทันตแพทย์และบุคลากรทางการแพทย์ (Dentist / Medical Personnel):} เป็นผู้ใช้งานหลักของระบบ (Primary Actor) ที่เข้ามาปฏิบัติงานทางคลินิก ได้แก่ การเข้าสู่ระบบ การเลือกสถาปัตยกรรมแบบจำลอง AI ที่ต้องการ การอัปโหลดภาพถ่ายรอยโรค การระบุอาการทางคลินิกและพฤติกรรมเสี่ยงของผู้ป่วย ตลอดจนการตรวจสอบผลการประเมินความเสี่ยงและภาพอธิบายผลทางคลินิก
    \item \textbf{ผู้ดูแลระบบ (System Administrator / Admin):} เป็นกลุ่มผู้ใช้งานที่เตรียมการรองรับสำหรับการขยายขอบเขตระบบในอนาคต (Future Administrative Actor) ทำหน้าที่บริหารจัดการระบบในภาพรวม ได้แก่ การจัดการบัญชีผู้ใช้งานและสิทธิ์การเข้าถึง การจัดการและตรวจสอบสถานะแบบจำลองปัญญาประดิษฐ์ ตลอดจนการตรวจสอบบันทึกประวัติการทำงานของระบบ (System \& Inference Logs)
\end{enumerate}

กรณีการใช้งานหลักของระบบ (Use Cases) จัดหมวดหมู่ตามขอบเขตหน้าที่การทำงานได้ดังนี้:
\begin{enumerate}
    \item \textbf{กลุ่มการยืนยันตัวตน (Authentication):}
    \begin{itemize}
        \item \textbf{UC-AUTH (Login / Authenticate):} การเข้าสู่ระบบเพื่อยืนยันตัวตนสำหรับทันตแพทย์และผู้ดูแลระบบ
    \end{itemize}
    \item \textbf{กลุ่มฟังก์ชันการตรวจคัดกรองทางคลินิกสำหรับทันตแพทย์ (Clinical Screening Functions):}
    \begin{itemize}
        \item \textbf{UC-01 (Select AI Model Architecture):} การเลือกสถาปัตยกรรมแบบจำลองปัญญาประดิษฐ์ที่ต้องการใช้งาน (เช่น แบบจำลองตรวจจับวัตถุ YOLO หรือแบบจำลองหลายโมดาลิตี Multimodal DenseNet-BERT)
        \item \textbf{UC-02 (Upload Lesion Image):} การอัปโหลดภาพถ่ายรอยโรคในช่องปากเข้าสู่ระบบเพื่อเตรียมการประมวลผล
        \item \textbf{UC-03 (Input Patient Symptoms):} การระบุข้อมูลอาการทางคลินิกและประวัติพฤติกรรมเสี่ยงของผู้ป่วย ทั้งการเลือกจากรายการตัวเลือก (Checkboxes) และการกรอกข้อความบรรยายอาการเพิ่มเติม
        \item \textbf{UC-04 (View Clinical Assessment Result):} การเรียกดูผลการประเมินทางคลินิก ภาพผลลัพธ์พร้อมเครื่องมืออธิบายผล (Bounding Box สำหรับ YOLO หรือ Grad-CAM Heatmap สำหรับ Multimodal) ค่าคะแนนความเชื่อมั่น และระดับความเสี่ยงของรอยโรค
    \end{itemize}
    \item \textbf{กลุ่มฟังก์ชันการบริหารจัดการระบบในอนาคตสำหรับผู้ดูแลระบบ (Administrative Functions - Future Extension):}
    \begin{itemize}
        \item \textbf{UC-ADM-01 (Manage User Accounts):} การจัดการข้อมูลบัญชีผู้ใช้งาน กำหนดสิทธิ์ และการอนุมัติการเข้าใช้งานระบบของทันตแพทย์
        \item \textbf{UC-ADM-02 (Manage AI Models):} การจัดการแบบจำลองปัญญาประดิษฐ์ การตรวจสอบสถานะการเชื่อมต่อ และการอัปเดตไฟล์น้ำหนักแบบจำลอง (Model Weights)
        \item \textbf{UC-ADM-03 (View System Logs):} การตรวจสอบบันทึกประวัติการทำงานของระบบ สถิติการเรียกใช้งานการอนุมานผล (Inference Logs) และข้อผิดพลาดที่เกิดขึ้น
    \end{itemize}
\end{enumerate}

ผังกรณีการใช้งานของระบบสนับสนุนการวินิจฉัยรอยโรคทางช่องปาก แสดงความสัมพันธ์ระหว่างตัวแสดง (Actors) และกรณีการใช้งานทั้งหมด ดังแสดงในรูปที่ \ref{fig:usecase_diagram}:

% แผนภาพ Use Case Diagram (เว้นกรอบ Placeholder สำหรับใส่ภาพกราฟิก)
\begin{figure}[htbp]
\centering
\fbox{
    \parbox{0.92\textwidth}{
        \centering
        \vspace{2.5cm}
        \textbf{[พื้นที่สำหรับใส่ภาพ: แผนภาพกรณีการใช้งานของระบบ (Use Case Diagram)]}\\[0.3cm]
        {\footnotesize \textit{(Use Case Diagram: Actor ทันตแพทย์ (Dentist), ผู้ดูแลระบบในอนาคต (Admin) และความสัมพันธ์กับ Use Cases)}}\\[0.2cm]
        {\scriptsize สามารถใส่ภาพ \texttt{images/usecase\_diagram.png} ลงในจุดนี้}
        \vspace{2.5cm}
    }
}
\caption{แผนภาพกรณีการใช้งานของระบบสนับสนุนการวินิจฉัยรอยโรคทางช่องปาก (Use Case Diagram)}
\label{fig:usecase_diagram}
\end{figure}

\subsection{ผังกิจกรรมการทำงานของระบบ (Activity Diagram)}
เพื่ออธิบายลำดับขั้นตอนการประมวลผลคำขอส่งตรวจตั้งแต่จุดเริ่มต้นจนถึงการแสดงผลลัพธ์ ผังกิจกรรมการทำงาน (Activity Diagram) ของระบบแสดงกระบวนการทำงานแบบเป็นขั้นตอน ดังแสดงในรูปที่ \ref{fig:activity_diagram}:

% แผนภาพ Activity Diagram (เว้นกรอบ Placeholder สำหรับใส่ภาพกราฟิก)
\begin{figure}[htbp]
\centering
\fbox{
    \parbox{0.92\textwidth}{
        \centering
        \vspace{2.5cm}
        \textbf{[พื้นที่สำหรับใส่ภาพ: แผนภาพกิจกรรมการทำงานของการประเมินรอยโรค (Activity Diagram)]}\\[0.3cm]
        {\footnotesize \textit{(Activity Diagram: Login $\rightarrow$ Select Model $\rightarrow$ Upload Image $\rightarrow$ Select Symptoms $\rightarrow$ Require NLP? $\rightarrow$ Submit $\rightarrow$ Route to Model $\rightarrow$ View Result)}}\\[0.2cm]
        {\scriptsize สามารถใส่ภาพ \texttt{images/activity\_diagram.png} ลงในจุดนี้}
        \vspace{2.5cm}
    }
}
\caption{แผนภาพกิจกรรมการทำงานของการประเมินรอยโรคและการเลือกแบบจำลอง (Activity Diagram)}
\label{fig:activity_diagram}
\end{figure}

ขั้นตอนการทำงานใน Activity Diagram สามารถอธิบายตามลำดับการไหลของกระบวนการได้ดังนี้:
\begin{enumerate}
    \item \textbf{การเข้าสู่ระบบ (Login):} ผู้ใช้งานทำการเข้าสู่ระบบยืนยันตัวตนเพื่อเข้าถึงหน้าการตรวจคัดกรอง
    \item \textbf{การเลือกแบบจำลอง (Select Model):} ผู้ใช้งานเลือกสถาปัตยกรรมแบบจำลองปัญญาประดิษฐ์ที่ต้องการใช้ในการตรวจรอยโรค
    \item \textbf{การอัปโหลดรูปภาพ (Upload Image):} อัปโหลดภาพถ่ายรอยโรคในช่องปากที่ต้องการประเมินเข้าสู่ระบบ
    \item \textbf{การเลือกอาการทางคลินิก (Select Symptoms Checkbox):} ติ๊กเลือกอาการและพฤติกรรมเสี่ยงจากรายการตัวเลือกที่กำหนด เช่น พฤติกรรมการสูบบุหรี่ ดื่มสุรา เคี้ยวหมาก อาการแสบร้อน หรือลักษณะรอยโรคที่สังเกตได้
    \item \textbf{การตรวจสอบเงื่อนไขความต้องการภาษาธรรมชาติ (Decision Node: Require NLP?):}
    \begin{itemize}
        \item \textbf{กรณีใช่ (YES):} หากแบบจำลองที่เลือกเป็นแบบจำลองในกลุ่มที่ต้องใช้การวิเคราะห์ภาษาธรรมชาติเชิงลึก (เช่น Multimodal WangchanBERTa) ระบบจะแสดงช่องให้ผู้ใช้กรอกข้อความบรรยายอาการและบริบทเพิ่มเติม (\texttt{Fill Symptoms Context}) จากนั้นจึงส่งข้อมูลต่อไปยังขั้นตอนส่งตรวจ
        \item \textbf{กรณีไม่ใช่ (NO):} หากแบบจำลองที่เลือกประมวลผลเฉพาะรูปภาพหรือใช้กฎเกณฑ์จับคู่แบบพื้นฐาน (เช่น YOLO) ผู้ใช้งานสามารถข้ามไปยังขั้นตอนการกดส่งตรวจ (\texttt{Submit}) ได้ทันที
    \end{itemize}
    \item \textbf{การกระจายเส้นทางการประมวลผล (Route based on selection):} Backend API Gateway จะทำการส่งคำขอและข้อมูลไปยัง Inference Pipeline ของแบบจำลองที่ผู้ใช้เลือกไว้ภายใน AI Inference Engine
    \item \textbf{การแสดงผลการประเมิน (View Result):} ระบบรวบรวมผลลัพธ์ คำนวณระดับความเสี่ยง และส่งคืนข้อมูลภาพพร้อมป้ายกำกับเพื่อแสดงผลแก่ผู้ใช้งานบนแดชบอร์ด
\end{enumerate}

% ------------------------------------------------------------------------------
% 3.4 การพัฒนาส่วนติดต่อผู้ใช้งาน
% ------------------------------------------------------------------------------
\section{การพัฒนาส่วนติดต่อผู้ใช้งาน (User Interface Design)}
\label{sec:ui_design}

\subsection{การออกแบบส่วนติดต่อผู้ใช้งานสำหรับการส่งตรวจ (Screening Interface)}
ส่วนติดต่อผู้ใช้ถูกพัฒนาขึ้นโดยเน้นความสะดวกและรวดเร็วในการปฏิบัติงานของทันตแพทย์ ประกอบด้วยฟังก์ชันการทำงานหลัก:
\begin{itemize}
    \item \textbf{ช่องทางเลือกและอัปโหลดภาพ:} รองรับการลากวางไฟล์ภาพ (Drag-and-Drop) พร้อมระบบแสดงตัวอย่างภาพ (Image Preview)
    \item \textbf{เมนูเลือกสถาปัตยกรรมแบบจำลอง (Model Selection Toggle):} ออกแบบเป็นสวิตช์เลือกที่ชัดเจน ให้แพทย์สามารถสลับระหว่างการตรวจจับตำแหน่งด้วย YOLO หรือการวิเคราะห์เชิงลึกด้วย Multimodal DenseNet-BERT ได้อย่างอิสระ
    \item \textbf{ฟอร์มระบุอาการและพฤติกรรมเสี่ยง:} แบ่งหมวดหมู่ชัดเจน ได้แก่ พฤติกรรมสูบบุหรี่ ดื่มสุรา เคี้ยวหมาก อาการแสบร้อนเมื่อกินเผ็ด แผลเจ็บปวด คราบขาวที่เช็ดลอกได้ และตำแหน่งที่พบรอยโรค
\end{itemize}

\subsection{การออกแบบแดชบอร์ดแสดงผลและการเปรียบเทียบผลลัพธ์}
หน้าต่างแสดงผลลัพธ์ทางคลินิก (Clinical Assessment Result View) ได้รับการออกแบบให้แสดงข้อมูลสำคัญครบถ้วนในหน้าจอเดียว:
\begin{enumerate}
    \item \textbf{การแสดงผลภาพเปรียบเทียบ:} แสดงภาพต้นฉบับคู่กับภาพผลลัพธ์ โดยหากเลือกโมเดล YOLO จะแสดงภาพที่มีกรอบ Bounding Box พร้อมค่าความมั่นใจรายกล่อง แต่หากเลือก DenseNet-BERT จะแสดงภาพความร้อน Grad-CAM Heatmap ชี้จุดที่โมเดลใช้ในการตัดสินใจ
    \item \textbf{การแสดงป้ายระดับความเสี่ยง (Risk Level Badge):} แสดงแถบสีเตือนระดับความเสี่ยงตามเกณฑ์ทางคลินิก เช่น สีแดงสำหรับความเสี่ยงสูง (High Risk), สีส้มสำหรับความเสี่ยงปานกลาง (Medium Risk), และสีเขียวสำหรับความเสี่ยงต่ำ (Low Risk)
\end{enumerate}

% ------------------------------------------------------------------------------
% 3.5 การจัดเตรียมสภาพแวดล้อม การคอนเทนเนอร์ไรซ์ และการทดสอบระบบอัตโนมัติ
% ------------------------------------------------------------------------------
\section{การจัดเตรียมสภาพแวดล้อม การคอนเทนเนอร์ไรซ์ และการทดสอบระบบอัตโนมัติ}
\label{sec:devops_and_testing}

\subsection{การจัดสภาพแวดล้อมด้วย Docker และ Docker Compose}
ระบบถูกเขียนไฟล์คอนฟิกูเรชัน \texttt{docker-compose.yml} เพื่อควบคุมการทำงานร่วมกันของทั้ง 3 เซอร์วิสอย่างเป็นเอกภาพ:
\begin{itemize}
    \item \textbf{คอนเทนเนอร์ Frontend (\texttt{oralscan-frontend}):} รันบนสภาพแวดล้อม Node.js ผ่าน Vite Development Server ให้บริการส่วนติดต่อผู้ใช้ที่พอร์ต 5173 โดยเชื่อมต่อไปยัง Backend ผ่านตัวแปรสภาพแวดล้อม \texttt{VITE\_API\_URL}
    \item \textbf{คอนเทนเนอร์ Backend Gateway (\texttt{oralscan-backend}):} ใช้ Base Image \texttt{python:3.11-slim} ติดตั้งแพ็กเกจ FastAPI, Uvicorn, HTTPX และ Pydantic ให้บริการที่พอร์ต 8000 และเชื่อมต่อไปยัง AI Engine ผ่านตัวแปรสภาพแวดล้อม \texttt{AI\_ENGINE\_URL}
    \item \textbf{คอนเทนเนอร์ AI Inference Engine (\texttt{oralscan-ai-engine}):} ใช้ Base Image ที่ติดตั้งไลบรารีสำหรับการคำนวณเชิงลึก ได้แก่ PyTorch, Torchvision, Transformers, OpenCV-headless, และ Ultralytics ให้บริการที่พอร์ต 8001
    \item \textbf{การเชื่อมต่อเครือข่ายภายใน (Bridge Network - \texttt{oralscan-network}):} ทั้งสามคอนเทนเนอร์เชื่อมต่อสื่อสารกันผ่านวงเครือข่ายเสมือนภายใน ช่วยเพิ่มความปลอดภัยเนื่องจาก AI Engine ไม่จำเป็นต้องเปิดพอร์ตออกสู่เครือข่ายอินเทอร์เน็ตภายนอกโดยตรง
\end{itemize}

\subsection{การทดสอบระบบอัตโนมัติ (Automated API Testing)}
คณะผู้จัดทำได้เขียนสคริปต์ทดสอบระบบอัตโนมัติ (\texttt{test\_api.py}) เพื่อตรวจสอบความถูกต้องของระบบก่อนการส่งมอบ ครอบคลุมกรณีทดสอบหลักดังแสดงในตารางที่ \ref{tab:test_cases}:

\begin{table}[htbp]
\centering
\small
\caption{ชุดกรณีทดสอบอัตโนมัติ (Automated Test Cases) สำหรับตรวจสอบการทำงานของระบบ}
\label{tab:test_cases}
\begin{tabular}{p{2.2cm}p{4.2cm}p{3.5cm}p{3.5cm}}
\toprule
\textbf{รหัสทดสอบ} & \textbf{เงื่อนไขและข้อมูลนำเข้า} & \textbf{ผลลัพธ์ที่คาดหวัง} & \textbf{เกณฑ์การผ่าน} \\
\midrule
TC-SYS-01 & ส่งคำขอ \texttt{GET /api/health} & สถานะ HTTP 200 ระบบพร้อมทำงาน & ได้รับ JSON \texttt{status: healthy} \\
TC-SYS-02 & ส่งภาพและอาการด้วยโมเดล \texttt{YOLO} & ได้รับผลการทำนายและ Bounding Box & ได้รับ Bounding Box ถูกต้องตาม Schema \\
TC-SYS-03 & ส่งภาพและอาการด้วยโมเดล \texttt{DenseNet-BERT} & ได้รับผลความเสี่ยงและภาพ Grad-CAM & ได้รับผลการวิเคราะห์และ Heatmap URL \\
TC-SYS-04 & ส่งไฟล์ที่ไม่ใช่ภาพถ่าย (เช่น .pdf, .txt) & สถานะ HTTP 400 Bad Request & ปฏิเสธคำขอและแจ้งเตือนรูปแบบไฟล์ \\
TC-SYS-05 & ส่งคำขอโดยระบุ \texttt{modelType} ผิดพลาด & สถานะ HTTP 400 Bad Request & ปฏิเสธคำขอและแนะนำโมเดลที่ถูกต้อง \\
TC-SYS-06 & ทดสอบการส่งคำขอพร้อมกัน 20 คำขอ & ทุกคำขอประมวลผลสำเร็จโดยไม่เกิด Timeout & อัตราความสำเร็จร้อยละ 100 \\
\bottomrule
\end{tabular}
\end{table}

% ------------------------------------------------------------------------------
% 3.6 สรุปวิธีการดำเนินงาน
% ------------------------------------------------------------------------------
\section{สรุปวิธีการดำเนินงาน}
\label{sec:methodology_summary}

การดำเนินงานในบทที่ 3 ได้มุ่งเน้นการแก้ปัญหาข้อจำกัดของงานวิจัยต้นแบบในอดีต โดยนำแบบจำลองตรวจจับวัตถุ YOLO จากรุ่นพี่กลุ่มที่ 1 และแบบจำลองการเรียนรู้เชิงลึกแบบหลายโมดาลิตีจากรุ่นพี่กลุ่มที่ 2 มารวบรวมและบูรณาการเข้าด้วยกันภายใต้สถาปัตยกรรมไมโครเซอร์วิส 3 คอนเทนเนอร์ (Frontend, API Gateway และ AI Engine) ที่ทันสมัยและประหยัดทรัพยากร พร้อมทั้งสร้างทางเลือกให้ทันตแพทย์สามารถเลือกใช้แบบจำลองได้ตามความต้องการของเคสผู้ป่วย (Alternative Model Selection) ผ่านการวางผังกรณีการใช้งาน (Use Case Diagram) และผังกิจกรรมการประมวลผลคำขอ (Activity Diagram) 

ผลการทดสอบการทำงานของระบบ ประสิทธิภาพความเร็วในการประมวลผล และผลการเปรียบเทียบความแม่นยำของแบบจำลองแต่ละแบบบนแพลตฟอร์มนี้ จะถูกนำเสนออย่างละเอียดในบทที่ 4 ต่อไป
